\section{Backend Tests} \label{sec:backend_tests}
We adopted a layered testing approach to ensure high confidence in our system’s behavior. With some tests written across the solution, our goal was to validate that every function in the application performs exactly as intended — both in isolation and when integrated with other core sub-systems. This rigorous test coverage helps us identify regressions early, improve maintainability, and confidently refactor our codebase during agile iterations.

\subsubsection{Unit Tests}
For the domain logic residing in our Service Layer, we wrote extensive unit tests using Moq to isolate external system behaviors, such as database repositories and cross-cutting infrastructure concerns. These tests validate:

\begin{itemize}
    \item Strict domain validation rules (e.g., preventing empty, white-spaced, or overly long conversation tag names).
    \item Behavior under diverse edge cases (e.g., attempting to tag a conversation where the user is not an active participant, or manipulating tags that belong to other identities).
    \item Error routing integrity using explicit strongly-typed results via the \texttt{OneOf} library pattern.
\end{itemize}

To keep our test suites maintainable, clean, and fast, we structured them using XUnit's \texttt{IClassFixture<T>} interface. We established a reusable fixture class for each core service lifecycle, such as the \texttt{ConversationTagServiceFixture}.

\begin{lstlisting}[caption={ConversationTagServiceFixture instantiation}, label={lst:conversation-tag-fixture-init}]
public class ConversationTagServiceTests : IClassFixture<ConversationTagServiceFixture>
\end{lstlisting}

Implementing centralized class fixtures yields significant advantages:
\begin{itemize}
    \item \textbf{Decoupled Mock Orchestration:} Pre-configures interconnected dependency graphs, mock states, and multi-repository contexts (such as \texttt{IConversationTagRepository}, \texttt{IUserRepository}, and \texttt{IConversationRepository}) within a centralized setup constructor.
    \item \textbf{Transactional Context Management:} Simulates nested transactional boundaries through specialized manager wrappers (\texttt{ITransactionManager}), mirroring real production database behaviors.
    \item \textbf{State Fluidity Utilities:} Exposes intuitive helper states (e.g., \texttt{SetupTagNameExists(bool)}, \texttt{SetupUserIsParticipant(bool)}) that let individual test cases manipulate mock flows cleanly using declarative state shifts.
\end{itemize}

\begin{lstlisting}[caption={Unit test execution inside ConversationTagServiceTests}, label={lst:conversation-tag-unit-tests}]
[Fact]
public async Task CreateTag_ShouldReturnError_WhenTagNameAlreadyExists()
{
    // Arrange
    _fixture.SetupUserExists(true);
    _fixture.SetupTagNameExists(true);
    _fixture.SetupCreateTag(10);

    // Act
    var result = await _fixture.ConversationTagService.CreateTag(_fixture.UserId, "Duplicada", "#FFFFFF");

    // Assert
    Assert.True(result.IsT1); // Validates that a failure union state was reached
    Assert.IsType<ConversationTagError.TagNameAlreadyExists>(result.AsT1);
}
\end{lstlisting}

As illustrated in Listing~\ref{lst:conversation-tag-unit-tests}, our test execution relies on declarative setup triggers. By feeding a pre-existing tag scenario into the mock layers, the service execution triggers a failure branches inside the domain validation. Instead of throwing raw exceptions, the application uses discriminated unions (\texttt{OneOf}) to output functional result branches. 

The test structure easily intercepts this by asserting that the response is mapped to the failure index state (\texttt{IsT1}) and confirms the dynamic payload matches the specific architectural error class (\texttt{ConversationTagError.TagNameAlreadyExists}).